\chapter{Safe Cracker}

\section{Overview}
In this project, you will build a combination-lock game using a rotary encoder and a small speaker.
Spin the encoder dial to select digits 0--9, press the built-in encoder button to lock in each digit,
and listen for audio hints that get warmer as you approach the correct value. Match the full secret
sequence to ``open the safe!''

At the end of this project, your microcontroller will run a fully playable combination-lock game
entirely driven by sound and a single knob—no screen required!

\begin{figure}[H]
\centering
    \maybeincludegraphics[width=.85\linewidth]{project_3/success.jpg}
    \caption{The assembled Safe Cracker, ready for a code to be cracked}
\end{figure}

\begin{tcolorbox}[colback=blue!5!white,colframe=blue!40!black]
    \textbf{Learning goals:}
    \begin{itemize}
        \item Wire a KY-040 rotary encoder and read its turns and button presses with interrupt-driven
              library code
        \item Wire a small speaker and play tones of different pitches on it with PWM (Pulse Width
              Modulation)
        \item Use circular distance to turn ``how close'' into hot-and-cold audio feedback
        \item Organize a game with a randomized secret code, state variables, callbacks, and a main loop
        \item Read the serial console to debug your program
    \end{itemize}
\end{tcolorbox}

\pagebreak

\section{Components}

\begin{components}
    \item Seeed XIAO ESP32-C3 microcontroller
    \item KY-040 rotary encoder module (5-pin, includes onboard pull-up resistors)
    \item Treedix 1W 8\si{\ohm} mini speaker
    \item Breadboard
    \item Jumper wires
\end{components}

\section{Directions}

\subsubsection{Remove previous components}
Before beginning, remove any components from prior chapters including LEDs, buttons, and wires.
Leave the microcontroller seated in the standard position shown in Project 1. Refer back to
that chapter's seating instructions if you need help.

\subsection{Understanding the KY-040 Rotary Encoder}

A rotary encoder is a type of position sensor that converts the rotation of a knob into electrical
signals. Unlike a volume knob that stops at each end (an electronic part called a
\textbf{potentiometer}), an encoder has no physical stops---you can spin it continuously in either
direction. As it turns, you feel small clicks; each click sends a short pattern of pulses that the
code counts.

The KY-040 module has five pins:
\begin{itemize}
    \item \textbf{CLK} — Clock signal; pulses as the knob turns
    \item \textbf{DT} — Data signal; used together with CLK to determine direction
    \item \textbf{SW} — Switch; goes LOW when the built-in button is pressed
    \item \textbf{+} — Power supply (connect to 3.3V); some boards label this \textbf{VCC}
    \item \textbf{GND} — Ground
\end{itemize}

Hold the module with the knob facing you and the pins pointing down. The white labels printed on
the circuit board, called the \textbf{silkscreen}, read
\textbf{GND}, \textbf{+}, \textbf{SW}, \textbf{DT}, \textbf{CLK} from left to right.

\begin{tcolorbox}[colback=yellow!10!white,colframe=yellow!50!black]
    \textbf{Important:} Power the KY-040 from \textbf{3.3V}, not 5V. The encoder's signal pins
    connect directly to the ESP32-C3's GPIO pins, which are built for 3.3V. If the encoder were
    powered from 5V, it would send 5V signals into those pins, which could damage your
    microcontroller.
\end{tcolorbox}

\subsection{Creating the circuit}

The encoder and speaker share the breadboard with the microcontroller.

\subsubsection{Place the components}

Seat the encoder module's five-pin header along the front edge of the breadboard, knob facing
you, so the pins land in silkscreen order: \textbf{GND} in
\bbhole{encoder.header.ground}{A12}, \textbf{+} in
\bbhole{encoder.header.power}{A13}, \textbf{SW} in
\bbhole{encoder.header.sw}{A14}, \textbf{DT} in
\bbhole{encoder.header.dt}{A15}, and \textbf{CLK} in
\bbhole{encoder.header.clk}{A16}. Row A leaves the knob clear of other components, making it
easier to grip and turn. Insert the speaker's two leads into the bottom half of the
breadboard: the red (+) lead in \bbhole{speaker.lead.signal}{C20} and the black ($-$) lead in
\bbhole{speaker.lead.ground}{C22}.

\begin{figure}[H]
    \centering
    \maybeincludegraphics[width=.95\linewidth]{project_3/components_placed.pdf}
    \caption{The encoder and speaker are placed, before adding jumper wires.}
\end{figure}

\subsubsection{Wire the KY-040 encoder}
\begin{center}
\begin{tabular}{|l|l|l|}
    \hline
    \textbf{KY-040 pin} & \textbf{XIAO pin} & \textbf{Jumper holes} \\
    \hline
    \connectionrow{encoder-ground}{GND}{GND}{I2}{E12}
    \connectionrow{encoder-power}{+ (VCC)}{3V3}{I3}{E13}
    \connectionrow{encoder-sw}{SW}{GPIO 8}{I6}{E14}
    \connectionrow{encoder-dt}{DT}{GPIO 9}{I5}{E15}
    \connectionrow{encoder-clk}{CLK}{GPIO 10}{I4}{E16}
    \hline
\end{tabular}
\end{center}

\subsubsection{Wire the speaker}
\begin{center}
\begin{tabular}{|l|l|l|}
    \hline
    \textbf{Speaker lead} & \textbf{XIAO pin} & \textbf{Jumper holes} \\
    \hline
    \connectionrow{speaker-signal}{Signal (+)}{GPIO 20}{I7}{E20}
    \connectionrow{speaker-ground}{Ground ($-$)}{GND}{J2}{E22}
    \hline
\end{tabular}
\end{center}

\begin{figure}[H]
    \centering
    \maybeincludegraphics[width=.95\linewidth]{project_3/wiring.pdf}
    \caption{The encoder and speaker are fully wired to the XIAO ESP32-C3.}
\end{figure}

\begin{tcolorbox}[colback=blue!5!white,colframe=blue!40!black]
    \textbf{How the game works:}
    A random 3-digit secret code (each digit 0--9) is generated when the program starts and printed
    to the serial console. Spin the encoder knob to change the current dial value. As you get closer
    to the target digit, the tone played on each tick gets higher and shorter. When you are exactly
    on the target, you hear a sharp high click. Press the encoder button to lock in your digit.
    A correct digit advances you to the next position; a wrong digit plays a low buzz and you try
    again. Crack all three digits to trigger the victory fanfare and start a new game!
\end{tcolorbox}

\subsection{Programming the microcontroller}

Connect the USB cable to the microcontroller and open the IDE. Refer back to the \nameref{ide}
chapter for instructions if needed.

Click on the file named \texttt{project\_3\_safe\_cracker.py}. Read through the comments and
code to get a sense of how it works, then click the blue play button to start the game.

When the program starts, watch the serial console. It prints the secret code so you can check
that the game is working (have a friend run it if you want a real challenge):

\begin{verbatim}
New safe code: [3, 7, 1]
Dial 1/3: showing 5 (target 3, distance 2)
\end{verbatim}

\begin{tcolorbox}[colback=green!5!white,colframe=green!40!black]
    \textbf{Tip:} If the tones are silent, check that the speaker's signal wire is connected
    to GPIO 20 and that the ground wire goes to GND. If the dial counts in the wrong direction,
    change \texttt{reverse=True} to \texttt{reverse=False} in the code.
\end{tcolorbox}

\subsection{Understanding the Code}

Here is a walk through the key sections of the code.

\subsubsection{Configuration Constants}

At the top of the file, a configuration block defines all pin assignments and game parameters
in one place. This makes it easy to adjust the game without hunting through the code:

\begin{lstlisting}[language=Python, caption=Configuration constants]
CODE_LENGTH = 3
DIAL_MIN = 0
DIAL_MAX = 9

PIN_CLK = 10
PIN_DT = 9
PIN_SW = 8
PIN_SPEAKER = 20

TONE_EXACT  = (880, 80)   # high click when exactly on target
TONE_WARM   = (550, 50)   # warm mid tone when 1-2 away
TONE_COLD   = (220, 30)   # low cold tone when 3+ away
TONE_BUZZ   = (150, 200)  # wrong-digit buzz
VICTORY_NOTES = [(262, 150), (330, 150), (392, 150), (523, 400)]
\end{lstlisting}

\subsubsection{Playing Tones with PWM}

The \texttt{play\_tone} function uses PWM (Pulse Width Modulation) to drive the speaker.
PWM rapidly switches the GPIO pin on and off at a given frequency, which the speaker cone
turns into sound:

\begin{lstlisting}[language=Python, caption=PWM tone helper]
def play_tone(freq, duration_ms):
    pwm = PWM(Pin(PIN_SPEAKER), freq=freq, duty=512)
    time.sleep_ms(duration_ms)
    pwm.deinit()
\end{lstlisting}

A \texttt{duty} of 512 out of 1023 means the pin is HIGH roughly half the time—this is
called a 50\% duty cycle and produces the loudest sound for a square wave.

\subsubsection{Proximity Hints}

The game measures how close the current dial value is to the target digit using a
\emph{circular distance}. Because the dial wraps around (0 and 9 are neighbors), the code
takes the smaller of the straight-line distance and the wrap-around distance:

\begin{lstlisting}[language=Python, caption=Circular distance on a 0-9 dial]
def dial_distance(current, target):
    diff = abs(current - target)
    return min(diff, 10 - diff)
\end{lstlisting}

This distance is then mapped to a tone:

\begin{center}
\begin{tabular}{|l|l|}
    \hline
    \textbf{Distance to target} & \textbf{Sound} \\
    \hline
    0 (exact match)             & High click (880 Hz, 80 ms) \\
    1 or 2 away                 & Warm mid tone (550 Hz, 50 ms) \\
    3 or more away              & Low cold tone (220 Hz, 30 ms) \\
    \hline
\end{tabular}
\end{center}

\subsubsection{Reading the Encoder and Button}

A knob can be turned at any moment, even while the program is busy playing a tone. If the
program only checked the pins now and then, it could miss clicks. Instead, it uses
\textbf{hardware interrupts}: the microcontroller watches the pins itself and, the instant one
changes, pauses whatever it was doing to run a short function. A function you hand to another
piece of code so it can call it later, when something happens, is called a \textbf{callback}.

The code uses two library modules from the \texttt{lib/} folder:

\begin{itemize}
    \item \textbf{RotaryIRQ} — reads the encoder using hardware interrupts; configured to count
          0--9 and wrap from 9 back to 0 (and from 0 to 9), like a real combination-lock dial.
    \item \textbf{DebouncedButton} — reads the encoder's push-button and fires a callback
          only once per press, filtering out electrical noise (bouncing).
\end{itemize}

\begin{lstlisting}[language=Python, caption=Setting up the encoder and button]
knob = RotaryIRQ(
    pin_num_clk=PIN_CLK,
    pin_num_dt=PIN_DT,
    min_val=DIAL_MIN,
    max_val=DIAL_MAX,
    reverse=True,
    pull_up=True,
    range_mode=RotaryIRQ.RANGE_WRAP,
)
knob.add_listener(on_dial_change)
DebouncedButton(PIN_SW, on_enter)
\end{lstlisting}

\texttt{reverse=True} makes clockwise rotation increase the displayed number (adjust to
\texttt{False} if your knob spins the wrong way). \texttt{pull\_up=True} enables the
internal pull-up resistors so the CLK and DT pins read HIGH when idle. \texttt{on\_dial\_change}
and \texttt{on\_enter} are the callbacks: the library calls them for you whenever the knob turns or
the button is pressed.

\subsubsection{Game Loop}

Each time the knob moves, \texttt{on\_dial\_change} is called automatically by the library.
It reads the current encoder value, computes the distance to the target digit, and plays
the corresponding hint tone. When the encoder button is pressed, \texttt{on\_enter} checks
whether the current value matches the target:

\begin{itemize}
    \item \textbf{Correct:} A confirming click plays, the dial index advances. If all digits
          are correct, the victory fanfare plays and a new game begins.
    \item \textbf{Wrong:} A low buzz plays and the player tries again on the same digit.
\end{itemize}

\section{Review}
In this project, you combined hardware input (rotary encoder) and output (PWM speaker tones) to build
a complete interactive game. You learned:
\begin{itemize}
    \item How a rotary encoder converts physical rotation into digital signals
    \item How PWM can drive a speaker to play different pitches
    \item How to use interrupt-driven library code to read hardware without blocking the main loop
    \item How circular distance can create smooth ``hot and cold'' proximity feedback
    \item How to structure a game with state variables, callbacks, and a main loop
\end{itemize}

\section{Possible Extensions}
If you want to experiment further, try these ideas:
\begin{itemize}
    \item Increase \texttt{CODE\_LENGTH} to 4 or 5 digits for a harder challenge
    \item Add a countdown timer—if you don't crack the safe in time, the game resets
    \item Change the victory fanfare to play a song you recognize (look up the note frequencies!)
    \item Add a ``lives'' system: after 3 wrong guesses on the same digit, the game resets
    \item Print the elapsed time to the serial console when the safe is cracked
\end{itemize}
